\section{ANEXO B. Entorno virtual
Python}\label{anexo-b.-entorno-virtual-python}

El desarrollo de una librería Python destinada a integrar múltiples
herramientas de análisis de seguridad plantea, desde el inicio, un
problema de gestión de dependencias que no debe subestimarse. Las
herramientas que se integran en la solución propuesta, como Slither,
Mythril o Echidna, tienen requisitos de versión específicos y en
ocasiones incompatibles entre sí cuando se instalan en el entorno global
del sistema. Esta situación, conocida en el ecosistema Python como
dependency hell, puede provocar conflictos silenciosos difíciles de
depurar y comprometer la reproducibilidad del entorno de desarrollo.

Un entorno virtual (virtual environment) es un directorio aislado que
contiene una instalación independiente del intérprete Python junto con
sus propios paquetes y dependencias, sin interferir con el sistema
global ni con otros proyectos. Esta separación proporciona varias
ventajas fundamentales en el contexto del presente trabajo:

En primer lugar, garantiza el \textbf{aislamiento de dependencias}, de
forma que cada proyecto mantiene sus propias versiones de bibliotecas
sin afectar al resto del sistema. Esto resulta especialmente relevante
cuando diferentes herramientas de análisis requieren versiones distintas
de una misma dependencia transitiva.

En segundo lugar, favorece la \textbf{reproducibilidad del entorno de
desarrollo}. Todos los integrantes del equipo pueden trabajar
exactamente con las mismas versiones de todas las dependencias,
eliminando la variabilidad asociada a instalaciones manuales y
garantizando que los resultados obtenidos durante el desarrollo son
consistentes independientemente del sistema operativo o configuración
personal de cada desarrollador.

En tercer lugar, facilita el \textbf{ciclo de vida} del proyecto al
delimitar claramente qué paquetes pertenecen al proyecto y cuáles son
del sistema, simplificando tanto la distribución de la librería como su
posterior publicación en registros públicos como PyPI.

\subsection{ANEXO B.1. Tipos de entornos
virtuales}\label{anexo-b.1.-tipos-de-entornos-virtuales}

En el ecosistema Python existen varias alternativas para la gestión de
entornos virtuales, con distintos niveles de abstracción y
funcionalidad.

El módulo estándar \texttt{venv}, incluido en la biblioteca estándar
desde Python 3.3, permite crear entornos virtuales básicos mediante el
comando \texttt{python\ -m\ venv\ .venv}. Sin embargo, este enfoque no
incluye gestión de dependencias ni ficheros de bloqueo
(\emph{lockfiles}), por lo que debe complementarse con herramientas
adicionales como \texttt{pip} y \texttt{pip-tools}.

\texttt{virtualenv} es una alternativa anterior al módulo estándar, con
mayor compatibilidad con versiones antiguas de Python y algunas
funcionalidades adicionales, aunque en la práctica ha quedado desplazada
por las herramientas modernas.

\texttt{conda} ofrece un modelo más completo que combina gestión de
entornos con gestión de paquetes, incluyendo dependencias no Python. Es
habitual en entornos científicos y de análisis de datos, pero introduce
una complejidad y un tamaño innecesarios para un proyecto centrado en el
ecosistema Python puro.

Herramientas como \texttt{poetry} o \texttt{pipenv} representan un nivel
superior de abstracción, combinando la gestión de entornos virtuales con
la resolución de dependencias, la generación de ficheros de bloqueo y el
ciclo de publicación de paquetes. Su adopción en proyectos profesionales
se ha generalizado en los últimos años.

Finalmente, \texttt{uv} constituye la herramienta de última generación
en este espacio, combinando todas las funcionalidades anteriores en una
solución de rendimiento muy superior, como se detalla en la sección
siguiente.

\subsection{ANEXO B.2. Herramienta a usar:
uv}\label{anexo-b.2.-herramienta-a-usar-uv}

\texttt{uv} es un gestor de paquetes y proyectos Python de alto
rendimiento desarrollado por Astral, la empresa creadora del formateador
Ruff. Implementado en Rust, se presenta como una herramienta unificada
capaz de sustituir a \texttt{pip}, \texttt{pip-tools}, \texttt{pipx},
\texttt{poetry}, \texttt{pyenv}, \texttt{twine} y \texttt{virtualenv}
mediante una interfaz de línea de comandos coherente. Su desarrollo se
orienta tanto a la velocidad de ejecución como a la corrección en la
resolución de dependencias y a la facilidad de adopción en proyectos de
diferente escala y complejidad.

\subsubsection{ANEXO B.2.1. Ventajas de
uv}\label{anexo-b.2.1.-ventajas-de-uv}

La principal ventaja diferencial de \texttt{uv} respecto a las
alternativas existentes es su rendimiento. Según los propios
\emph{benchmarks} publicados en su documentación oficial, \texttt{uv} es
entre 10 y 100 veces más rápido que \texttt{pip} en operaciones de
instalación de paquetes con caché caliente. Esta diferencia resulta
perceptible en la práctica, especialmente durante las fases de
incorporación de nuevos integrantes al equipo o en entornos de
integración continua donde el entorno debe reconstruirse frecuentemente.

Más allá del rendimiento, \texttt{uv} ofrece un conjunto de ventajas
relevantes para el desarrollo de la librería propuesta en este trabajo.
Gestiona automáticamente los entornos virtuales asociados a cada
proyecto, sin necesidad de crearlos ni activarlos manualmente. Genera y
mantiene un fichero de bloqueo universal (\texttt{uv.lock}) que
garantiza instalaciones reproducibles en cualquier plataforma. Permite
gestionar múltiples versiones del intérprete Python e instalar la
versión adecuada de forma automática si no está disponible en el
sistema. Además, su diseño es compatible con los estándares del
ecosistema Python (\texttt{pyproject.toml}, PEP 517, PEP 621), lo que
facilita la integración con otras herramientas y la publicación en
registros de paquetes.

\subsubsection{ANEXO B.2.2. Qué proporciona uv en el contexto de este
proyecto}\label{anexo-b.2.2.-quuxe9-proporciona-uv-en-el-contexto-de-este-proyecto}

Para el desarrollo de la librería propuesta, uv proporciona un conjunto
de funcionalidades que cubren todo el ciclo de vida del proyecto, desde
la inicialización hasta la publicación.

\textbf{Gestión de dependencias y sincronización del entorno.} Una vez
definidas las dependencias del proyecto en el fichero
\texttt{pyproject.toml}, cualquier integrante del equipo puede
reproducir exactamente el mismo entorno ejecutando un único comando:

\begin{Shaded}
\begin{Highlighting}[]
\ExtensionTok{uv}\NormalTok{ sync}
\end{Highlighting}
\end{Shaded}

Este comando resuelve las dependencias declaradas, instala las versiones
exactas registradas en el fichero de bloqueo \texttt{uv.lock} y
configura el entorno virtual del proyecto de forma automática. La
simplicidad de este flujo elimina los problemas habituales de
divergencia entre entornos de desarrollo individuales, garantizando que
todos los miembros del equipo trabajan con las mismas versiones de
Slither, Mythril, Echidna y el resto de dependencias de la librería.

\textbf{Inicialización de proyectos.} \texttt{uv} proporciona soporte
integrado para la creación de nuevos proyectos mediante el comando
\texttt{uv\ init}. Para el caso específico de una librería Python
destinada a ser importada por otros proyectos o publicada en PyPI, se
utiliza la opción \texttt{-\/-lib}:

\begin{Shaded}
\begin{Highlighting}[]
\ExtensionTok{uv}\NormalTok{ init }\AttributeTok{{-}{-}lib}\NormalTok{ evmaudit}
\end{Highlighting}
\end{Shaded}

Este comando genera automáticamente la estructura de directorios
recomendada para una librería Python, incluyendo el fichero
\texttt{pyproject.toml} con los metadatos del proyecto, el directorio
\texttt{src/} con el paquete principal y los ficheros de configuración
necesarios para la construcción y distribución. El uso de la disposición
\texttt{src/} (\emph{src layout}) es la práctica recomendada actualmente
para proyectos publicables, ya que evita problemas habituales
relacionados con la importación del paquete desde el directorio raíz
durante el desarrollo.

\textbf{Construcción de distribuciones.} \texttt{uv} integra soporte
nativo para la generación de distribuciones instalables mediante el
comando \texttt{uv\ build}, que produce tanto el archivo fuente
(\emph{sdist}) como la rueda binaria (\emph{wheel}) del paquete:

\begin{Shaded}
\begin{Highlighting}[]
\ExtensionTok{uv}\NormalTok{ build}
\end{Highlighting}
\end{Shaded}

El resultado son los artefactos estándar de distribución Python ubicados
en el directorio \texttt{dist/}, listos para ser publicados o
distribuidos directamente.

\textbf{Publicación de paquetes.} El ciclo se completa con el soporte
para publicación en registros de paquetes, incluyendo PyPI, mediante el
comando \texttt{uv\ publish}:

\begin{Shaded}
\begin{Highlighting}[]
\ExtensionTok{uv}\NormalTok{ publish}
\end{Highlighting}
\end{Shaded}

Este comando gestiona la autenticación y la subida de los artefactos
generados, cubriendo el flujo completo que anteriormente requería
herramientas adicionales como \texttt{twine}.

En conjunto, \texttt{uv} unifica en una sola herramienta todo el ciclo
de vida del proyecto: inicialización, gestión de dependencias,
sincronización del entorno, construcción de distribuciones y
publicación. Esta integración reduce la fricción en el desarrollo
colaborativo y facilita la adopción de prácticas profesionales de
gestión de proyectos Python desde las primeras fases del trabajo.
